\section{Three-mode floor chambers and exact small-grid values}\label{sec:floor-chambers}
For three modes, strict cumulative-floor chambers have a particularly small
combinatorial description. It determines the exact minimax values on five,
six, and seven equal cells at the budgets considered below, closes a
continuous boundary case, and identifies a concrete obstruction to extending
the five-cell construction to arbitrary budgets.

\subsection{A complete three-state description}
Work on $N$ unit cells and suppose initially that $A_i(j)\notin\mathbb Z$
for every mode $i\in\{0,1,2\}$ and prefix $1\le j\le N$. Put
$f_j=(\lfloor A_i(j)\rfloor)_i$ and
$\sigma_j=j-\sum_i f_{j,i}$. Since all three fractional parts are strictly
between zero and one, $\sigma_j\in\{1,2\}$. An integer word has discrepancy
less than one exactly when its prefix counts lie coordinatewise between
$f_j$ and $f_j+\mathbf1$. Its possible count vectors at this prefix are
therefore exactly the three vectors
\begin{equation}\label{eq:floor-offset}
 f_j+o_{\sigma_j}(u),\quad u\in\{0,1,2\},\qquad
 o_1(u)=e_u,\quad o_2(u)=\mathbf1-e_u.
\end{equation}
The index $u$ labels a count vector; it need not be the active mode.

\begin{proposition}[Realizable floor histories]\label{prop:floor-automaton}
A sequence of integer floors is realized by a relaxed control with all
prefix coordinates nonintegral if and only if it starts with
$f_1=0$, $\sigma_1=1$ and, at each subsequent cell, its increment
$d=f_j-f_{j-1}$ has one of the following forms:
\begin{equation}\label{eq:floor-transitions}
\begin{array}{c|c|c}
\sigma_{j-1}&\sigma_j&d\\\hline
1&1&e_p\quad(p=0,1,2)\\
1&2&0\\
2&1&\mathbf1-e_p\quad(p=0,1,2)\\
2&2&e_p\quad(p=0,1,2).
\end{array}
\end{equation}
For each such transition there is an edge from count node $u$ to count
node $v$, labeled by active mode $a$, precisely when
\begin{equation}\label{eq:floor-edge}
 d+o_{\sigma_j}(v)-o_{\sigma_{j-1}}(u)=e_a.
\end{equation}
The full integer words in this layered graph are exactly the words with
discrepancy less than one throughout this floor chamber.
\end{proposition}
\begin{proof}
Every cell allocation lies in $[0,1]$, so each floor increment is zero or
one. Its sum is $1+\sigma_{j-1}-\sigma_j$, giving precisely
\eqref{eq:floor-transitions}; the initial floors vanish since three
positive nonintegral allocations sum to one at the first prefix.
An integer word increments its count vector by one unit vector per cell,
which proves \eqref{eq:floor-edge} and the word correspondence whenever
the chamber exists.

It remains to prove that every displayed history really exists. Each of
its transition graphs has an incoming and an outgoing edge at every node.
This follows directly from \eqref{eq:floor-edge}. More explicitly,
for $1\to1$ with increment $e_p$, all self edges have label $p$ and
$u=p$ can move to any $v$ with label $v$.
For $1\to2$, the edges are $u\ne v$, labeled by the third mode.
For $2\to1$ with increment $\mathbf1-e_p$, the edges are $u=p$,
labeled $v$, or $v=p$, labeled $u$.
For $2\to2$ with increment $e_p$, all self edges have label $p$ and
any $u$ can move to $v=p$ with label $u$.
All three initial nodes are available, with first active mode $u$.
Induction forward and backward shows that every node at every layer lies
on a full path. Average the cell allocations of all these finitely many
paths. At every prefix and coordinate, the three count nodes include both
the floor and ceiling, and all three occur on full paths. Their average
is therefore strictly between these endpoints. This average is a valid
simplex-valued relaxed control realizing the prescribed floors.
\end{proof}
This proof gives a realizability criterion without a linear program or a
search through infeasible floor arrays. The numbers of histories ending
in states $1,2$ are the row vector
\begin{equation}\label{eq:floor-count}
 (1,0)\begin{pmatrix}3&1\\3&3\end{pmatrix}^{N-1}.
\end{equation}
The total counts for $N=1,\ldots,7$ are
$1,4,18,84,396,1872,8856$.

The least switch count in one chamber is also computed by a small exact
recursion. Let $C_j(u,a)$ be the least count among its prefix paths ending
at node $u$ with last active mode $a$, and use $+\infty$ for an empty set.
Initialize $C_1(u,a)=0$ if $a=u$ and $+\infty$ otherwise. For each edge
$u\to v$ labeled $b$, update
\begin{equation}\label{eq:floor-cost-DP}
 C_j(v,b)=\min_{\substack{u,a:\,u\to v\text{ has label }b}}
       \{C_{j-1}(u,a)+\mathbf1_{a\ne b}\}.
\end{equation}
Every prefix word has exactly one predecessor, so induction proves this
recursion, including the unrestricted initial activation convention.
The answer is $\min_{u,a}C_N(u,a)$.

\subsection{The exact five-cell theorem and a continuous consequence}
\begin{theorem}[Five equal cells]\label{thm:five-cells}
For three modes and five cells of common length $\Delta>0$,
\begin{equation}\label{eq:five-cells}
 F^{\mathcal T}_{3,2}=\Delta.
\end{equation}
The same value holds when the relaxed input is any measurable control.
\end{theorem}
\begin{proof}
Scale to unit cells. Enumerate the four transition types in
\eqref{eq:floor-transitions} for four successive steps from
$(f_1,\sigma_1)=(0,1)$, taking all three choices of $p$ where present.
For every resulting history apply \eqref{eq:floor-edge} and
\eqref{eq:floor-cost-DP}. The exact integer computation gives
\begin{center}
\begin{tabular}{c|rrr|r}
minimum switches&0&1&2&total\\\hline
number of strict chambers&3&138&255&396
\end{tabular}
\end{center}
There are no other histories by Proposition~\ref{prop:floor-automaton},
and all transition branches are enumerated. The standard-library checker
\texttt{verification/stage04/check\_floor\_chambers.py} implements these
three displayed formulas and verifies each terminal minimum is at most two.
Thus this is a finite integer certificate, not a numerical estimate.
An independent bundled checker also enumerates all $3^5=243$ words against
all $1\cdot4\cdot9\cdot16\cdot25=14400$ candidate floor arrays and
recovers the same chambers and switch counts.
For that alternative enumeration, strict feasibility is equivalent to each
bounded prefix coordinate attaining both its floor and ceiling among the
feasible integer words. Necessity follows by decomposing a feasible
fractional prefix flow into integral flows; sufficiency follows by averaging
all the feasible words. This gives an independent coverage check rather
than relying only on agreement of the final count.

Every control with nonintegral prefixes consequently has a two-switch
word with error strictly below one. To cover integral boundary coordinates,
perturb the cell rates toward $(1/7,2/7,4/7)$. Every coordinate of this
constant control is nonintegral at prefixes $1,\ldots,5$. Along a sequence
of positive perturbations tending to zero all prefix coordinates avoid
integers: each coordinate has only finitely many forbidden perturbation
parameters unless constant, in which case it is already nonintegral.
The minimum discrepancy over the finite set of two-switch words is
continuous, so the original control has error at most one.
For arbitrary $N$ the same argument uses, for example,
$(1/p,2/p,(p-3)/p)$ with a prime $p>\max\{N,3\}$.

For sharpness, take the pure relaxed word $01010$. Its integer prefix
counts force any word with error less than one to have the same prefix
counts and hence the identical five letters, requiring four switches.
Thus every two-switch word has error at least one. Cell averaging and
scaling complete the proof.
\end{proof}

\begin{corollary}[Three modes, two continuous switches]\label{cor:three-two-continuous}
For every $T>0$,
\begin{equation}\label{eq:three-two-continuous}
 F_{3,2}(T)=T/5.
\end{equation}
\end{corollary}
\begin{proof}
The upper bound follows from \eqref{eq:minimax-grid-domination} and
Theorem~\ref{thm:five-cells} on five cells of length $T/5$.
The lower bound is the case $(n,s)=(3,2)$ of
\citet[Proposition~4, p.~612, equation~(7.4)]{sager-zeile2021}.
Here is a direct alternative witness. Scale to $T=5$ and take five pure
unit phases $01201$, with terminal masses $(2,2,1)$.
If a schedule had error $E<1$, it would need to use all three modes,
so with at most three blocks each mode would occur exactly once.
Let $i$ be its final mode, starting at $v$.
The first times when the three relaxed allocations reach one are
$d=(1,2,3)$. Since mode $i$ is unserved before $v$, necessarily
$v<d_i$. On the other hand its final signed error gives
$v\ge5-m_i-E>5-m_i-1$, whose values are $(2,2,3)$ and are
coordinatewise at least $d$. This is a contradiction.
The order $012$ with switches at $2,4$ has error exactly one,
as endpoint monotonicity or direct integration verifies.
\end{proof}
This adds the boundary $n=3$ to the continuous two-switch result of
Theorem~\ref{thm:small-full}. It concerns full absolute error and does
not assert uniform extremality for the one-sided three-block problem.

\begin{corollary}[Six equal cells]\label{cor:six-cells}
For three modes, six cells of length $\Delta$, and at most three switches,
$F^{\mathcal T}_{3,3}=\Delta$.
\end{corollary}
\begin{proof}
The same complete transition enumeration at depth six gives
$3,255,1377,237$ chambers with respective minimum switch counts
$0,1,2,3$, totaling $1872$. Thus every strict chamber admits a
three-switch word of error less than one after scaling to $\Delta=1$.
The preceding perturbation argument covers its boundary.
The pure relaxed word $010101$ gives the reverse bound: error below one
would force all five switches of that word.
\end{proof}

\subsection{The sharp seven-cell value}
\begin{proposition}[A three-switch grid instance]\label{prop:seven-cell-instance}
Let $u=(1/3,1/3,1/3)$ and take seven unit-cell rate vectors
\begin{equation}\label{eq:seven-cell-rates}
 (u,e_0,e_1,e_2,e_0,u,e_2).
\end{equation}
The optimal grid discrepancy with at most three switches is $4/3$.
Consequently $F^{\mathrm{unit}}_{3,3}(7)\ge4/3$.
\end{proposition}
\begin{proof}
The cumulative vectors, multiplied by three, are
\[
 (1,1,1),\ (4,1,1),\ (4,4,1),\ (4,4,4),\
 (7,4,4),\ (8,5,5),\ (8,5,8).
\]
Each coordinate has fractional part $1/3$ or $2/3$. An integer outside
its floor--ceiling pair therefore differs from it by at least $4/3$.
Any word of discrepancy less than $4/3$ must lie in this floor chamber.
For clarity, the entire recursion \eqref{eq:floor-cost-DP} for this chamber
is displayed below. Each parenthesized triple is one row of the matrix
$C_j(u,a)$, with rows indexed by $u$ and columns by $a$ in the order
$0,1,2$:
\[
\begin{array}{c|ccc}
 j&\text{row }0&\text{row }1&\text{row }2\\\hline
1&(0,\infty,\infty)&(\infty,0,\infty)&(\infty,\infty,0)\\
2&(0,\infty,\infty)&(1,1,\infty)&(1,\infty,1)\\
3&(1,1,\infty)&(\infty,1,\infty)&(\infty,2,2)\\
4&(3,\infty,2)&(\infty,2,2)&(\infty,\infty,2)\\
5&(3,\infty,\infty)&(3,3,\infty)&(3,\infty,2)\\
6&(\infty,3,4)&(3,\infty,4)&(3,4,\infty)\\
7&(\infty,\infty,4)&(\infty,\infty,4)&(4,4,4)
\end{array}
\]
The final minimum is four, proving the lower bound under three switches.
The word $0011022$ has three switches and direct evaluation of its integer
prefix counts gives maximum discrepancy $4/3$, proving attainment.
The checker independently enumerates all $3^7=2187$ words: exactly
$104$ belong to the floor chamber and their minimum switch count is four.
\end{proof}
The exceptional chamber in this example is sufficient to determine the
full seven-cell minimax value as well.

\begin{theorem}[Seven equal cells]\label{thm:seven-cells}
For three modes, seven cells of length $\Delta$, and at most three switches,
\begin{equation}\label{eq:seven-cells}
 F^{\mathcal T}_{3,3}=4\Delta/3.
\end{equation}
\end{theorem}
\begin{proof}
Scale to $\Delta=1$. Exhaustive evaluation of
\eqref{eq:floor-transitions}--\eqref{eq:floor-cost-DP} gives
\begin{center}
\begin{tabular}{c|rrrrr|r}
minimum switches&0&1&2&3&4&total\\\hline
number of strict chambers&3&414&4542&3891&6&8856
\end{tabular}
\end{center}
The six chambers requiring four switches are exactly the six mode
permutations of the floor history in
Proposition~\ref{prop:seven-cell-instance}:
\begin{equation}\label{eq:seven-bad-floor}
 (0,0,0),\ (1,0,0),\ (1,1,0),\ (1,1,1),\
 (2,1,1),\ (2,1,1),\ (2,1,2).
\end{equation}
The checker tests equality of these two sets of histories, rather than only
their cardinalities. Every other strict chamber has a three-switch word
of discrepancy less than one.

Consider \eqref{eq:seven-bad-floor}; relabeling handles its permutations.
The following words have at most three switches:
\[
 w^0=1200022,\qquad w^1=0021122,\qquad w^2=0011220.
\]
The prefix counts of $w^i$ belong to every required floor--ceiling pair
except coordinate $i$ at time $i+2$, where their count is zero instead
of one or two. This is a direct integer check over the displayed floor
history. Consequently, for any relaxed control in this chamber,
\[
 D(A,w^0)\le\max\{1,A_0(2)\},\quad
 D(A,w^1)\le\max\{1,A_1(3)\},\quad
 D(A,w^2)\le\max\{1,A_2(4)\}.
\]
Monotonicity and conservation give
$A_0(2)+A_1(3)+A_2(4)\le\sum_i A_i(4)=4$.
At least one of these three coordinates is therefore at most $4/3$,
and its corresponding word proves the desired upper bound.
Generic perturbation and continuity cover all boundary profiles.
Proposition~\ref{prop:seven-cell-instance} supplies the matching lower
bound, and cell averaging includes measurable inputs.
\end{proof}
Thus error one on $2s+1$ unit cells fails already at $s=3$, while the
sharp replacement there is $4/3$. This value is below the published
conjectured upper bound $3/2$ at the same parameters, so it is not a
counterexample to that conjecture. The continuous consequences are
\begin{equation}\label{eq:three-three-band}
 \frac T7\le F_{3,3}(T)\le\frac T6.
\end{equation}
The upper bound follows from Corollary~\ref{cor:six-cells} and
\eqref{eq:minimax-grid-domination}. The lower bound is
\citet[Proposition~4, p.~612]{sager-zeile2021}; it can also be checked
directly on that source's pure seven-unit-cell word $0120011$.
If $E<1$, modes $0,1,2$ must occur before times $1,2,3$, respectively.
Mode $0$ must also occur in $(3,5)$, and mode $1$ in $(5,7)$, since
otherwise the corresponding error changes by two across that interval.
A single activation of mode $0$ spanning its early and late occurrences
would cover $[1,3]$, where its relaxed rate is zero, again changing
its error by two. A single activation of mode $1$ would cover $[2,5]$,
changing its error by three. Both changes exceed $2E$.
Thus modes $0$ and $1$ each need two activations and mode $2$ needs one,
requiring at least five blocks. Scaling proves the lower bound.
This interval does not identify an
exact continuous three-switch value. In particular, the sharp seven-cell
value is a grid theorem, and its transfer would give only the weaker
upper bound $4T/21$.

\subsection{The scope of the published lower-bound correction}
\citet[Corollary~5, p.~611]{sager-zeile2021} state, for
$n_\omega>2$ and $1\le\sigma_{\max}\le N-2$, the lower bound
\begin{equation}\label{eq:source-cor5}
 \theta^{\max}\ge
 \frac{N+\sigma_{\max}+1}{3+2\sigma_{\max}}\,\bar\Delta.
\end{equation}
At $(n_\omega,N,\sigma_{\max})=(3,5,2)$ this is $8\Delta/7$,
whereas Theorem~\ref{thm:five-cells} gives exactly $\Delta$.
All the stated parameter restrictions hold. Even if the source's outer
maximum is read as varying the grid with fixed horizon $5\Delta$, five
cells, and maximum width $\Delta$, every cell must equal $\Delta$.
Thus that reading does not remove the contradiction.

The preceding attainment argument in Corollary~4, p.~610, uses
$N=j(3+2\sigma_{\max})+\sigma_{\max}+2$ with integer $j\ge0$.
It does not establish attainment at every $N$; this restriction is absent
from Corollary~5. The five-cell theorem disproves
\eqref{eq:source-cor5} as written. This correction is separate from the
many-mode counterexamples to Conjecture~1; it supplies no replacement
lower formula for the full parameter range.
